\chapter{Introduction}

\section{Motivation}

Software vulnerabilities are a persistent and growing challenge in computer security.
The National Vulnerability Database recorded over 28{,}000 new CVE entries in 2023 alone,
spanning memory safety errors, injection flaws, cryptographic weaknesses, and access control failures
across virtually every class of software system~\citep{cvedetails2023}.
The consequences of undetected vulnerabilities range from data breaches and service outages
to compromised infrastructure, with the economic cost of software vulnerabilities estimated
in the hundreds of billions of dollars annually~\citep{mcafee2020cybercrime}.

Despite this scale, the primary method by which organisations detect vulnerabilities before
deployment remains manual code review --- a process that is slow, expensive, and dependent on
the reviewer's familiarity with the specific vulnerability class and codebase.
Automated static analysis tools such as Flawfinder and Coverity partially automate this process,
but they operate on hand-written rule sets that require expert maintenance, produce high rates of
false positives that erode analyst trust, and cannot generalise to vulnerability patterns
not anticipated by their authors.

Machine learning approaches offer a different trade-off: rather than encoding expert rules,
they learn vulnerability patterns from labelled examples, with the potential to discover
surface regularities that human-written rules miss.
Recent years have seen substantial progress in applying deep learning to this problem,
from LSTM-based approaches operating on token sequences to graph neural networks that
incorporate program structure, to transformer encoders pre-trained on hundreds of millions
of lines of code.
Yet two fundamental problems remain largely unsolved.
First, existing systems are predominantly language-specific: a model trained on C vulnerability
data is not directly applicable to Python, and vice versa, even when the underlying flaw
is conceptually identical.
Second, models trained on benchmark datasets of real-world labelled code often fail to
generalise in practice, either because the benchmark distribution does not reflect the
diversity of production code, or because the model has learned surface correlates of
vulnerability labels rather than the underlying vulnerability patterns.

\section{Research Questions}

This dissertation addresses these problems through a set of concrete research questions:

\begin{enumerate}
  \item \textbf{RQ1}: Can a pre-trained code model fine-tuned on a synthetic multi-language
        dataset reliably detect vulnerabilities across C and Python?
  \item \textbf{RQ2}: How does joint cross-language training affect per-language and
        cross-language detection performance relative to single-language models?
  \item \textbf{RQ3}: To what extent does a model trained on synthetic data generalise
        to real-world vulnerability patterns, and what is the magnitude of the performance gap?
  \item \textbf{RQ4}: Do targeted training interventions --- hard negative mining and
        real-world data augmentation --- measurably reduce the synthetic-to-real gap?
  \item \textbf{RQ5}: Is the model performing keyword matching or genuine semantic
        pattern recognition, and which failure modes are architectural rather than correctable
        through data?
\end{enumerate}

\section{Approach and Contributions}

The system developed in this dissertation fine-tunes UniXcoder~\citep{guo2022unixcoder},
a 125-million-parameter encoder pre-trained on code and natural language across multiple
programming languages, as a binary vulnerability classifier.
A lightweight classification head --- a dropout layer followed by a linear projection
and sigmoid activation over the \texttt{[CLS]} token embedding --- is attached to
the pre-trained encoder and trained end-to-end on a hand-crafted multi-language
synthetic dataset.

The synthetic dataset covers 17 CWE types in C and Python, with 8 shared cross-language
vulnerability classes, and is constructed to include clean contrastive pairs:
each vulnerable example has a corresponding safe counterpart that uses the correct
mitigation while preserving the surrounding code structure.

Four quantitative experiments evaluate the system: a baseline performance evaluation on
a held-out synthetic set (Experiment 1), a cross-language ablation comparing C-only,
Python-only, and joint models in a 3$\times$3 matrix (Experiment 2), a real-world
generalisation test on 200 DiverseVul~\citep{ni2023diversevul} C functions
(Experiment 3), and a decision threshold sweep (Experiment 4).
A qualitative semantic understanding test suite with four probing tests distinguishes
pattern-matching from structural reasoning.

Two interventions are implemented and evaluated: hard negative mining to resolve
false positives on safe strong-cryptography functions, and real-world data augmentation
with 1{,}000 DiverseVul training samples to reduce the synthetic-to-real gap.

The main contributions are:
\begin{itemize}
  \item A parameterised multi-language synthetic vulnerability dataset covering 17 CWE types,
        with contextual wrapping and hard negative contrastive pairs.
  \item A fine-tuned UniXcoder classifier achieving F1\,=\,0.931 on synthetic evaluation
        and F1\,=\,0.749 on real-world evaluation.
  \item A cross-language ablation study quantifying transfer quality at the CWE level and
        identifying the structural conditions for successful transfer.
  \item A semantic understanding test suite with four probing tests that locate the model's
        architectural ceiling precisely.
  \item Token-level attribution analysis via occlusion, quantifying which tokens drive each
        prediction and confirming the mechanistic basis of the model's successes and failures.
  \item Empirical evaluation of hard negative mining and real-world augmentation as targeted
        interventions, with before-and-after comparison.
\end{itemize}

\section{Dissertation Structure}

Chapter~\ref{ch:background} reviews related work, covering static analysis tools,
machine learning approaches to vulnerability detection, pre-trained code models,
vulnerability datasets, and calibration methods.
Chapter~\ref{ch:methodology} describes the model architecture, dataset construction,
training procedure, and evaluation framework.
Chapter~\ref{ch:experiments} presents results from all four experiments and the
semantic understanding tests.
Chapter~\ref{ch:analysis} analyses cross-language transfer, real-world error patterns,
calibration, the per-CWE generalisation gap, and token-level attribution.
Chapter~\ref{ch:improvements} describes the two training interventions and quantifies
their effect.
Chapter~\ref{ch:discussion} discusses limitations and future directions.
Chapter~\ref{ch:conclusion} concludes.
